\section{Quản trị tài chính và Nguồn lực}

\subsection{Dự toán tài chính chi tiết}
Hệ thống tài chính của CourtMate được xây dựng để hỗ trợ mô hình kinh doanh B2B SaaS (Phần mềm dịch vụ dành cho doanh nghiệp), tập trung vào việc tối ưu hóa dòng tiền và duy trì lợi thế cạnh tranh "Zero-Commission".
\subsubsection{Vốn đầu tư ban đầu (CAPEX)}
Đây là nguồn vốn tập trung vào việc xây dựng nền tảng từ đầu đến khi sẵn sàng tung ra thị trường (giai đoạn MVP)
\begin{itemize}
    \item Chi phí Nghiên cứu và Phát triển (R\&D): dao động khoảng 150 - 200 triệu VNĐ. Khoản này chi trả cho việc thiết kế kiến trúc hệ thống, xây dựng các quy trình nghiệp vụ đặt sân và đảm bảo tính bảo mật dữ liệu.
    \item Thiết lập hạ tầng máy chủ: triển khai mô hình Multi-Cloud để tận dụng tốc độ của nhà cung cấp trong nước và sự ổn định của hệ thống quốc tế.
\end{itemize}

\subsubsection{Chi phí vận hành hàng tháng (OPEX)}
\noindent \textbf{Lựa chọn nhà cung cấp}: nhóm thực hiện đối sánh năng lực giữa các đơn vị cung cấp máy chủ lớn hiện nay
\begin{table}[htbp]
    \centering
    \begin{tabularx}{\textwidth}{@{}l X X X@{}}
        \toprule
        \textbf{Tiêu chí} & \textbf{AWS (Toàn cầu)} & \textbf{Viettel IDC (Nội địa)} & \textbf{Bizfly Cloud (Nội địa)} \\
        \midrule
        Độ trễ & Cao (30-50ms) & Thấp (<10ms) & Thấp (<10ms) \\
        Chi phí tính toán & Khá cao & Trung bình & Thấp (tối ưu nhất) \\
        Quản lý dữ liệu & Ổn định & Đang phát triển & Đang phát triển \\
        Chi phí băng thông & Cao & Miễn phí/Rẻ & Miễn phí/Rẻ \\
        \bottomrule
    \end{tabularx}
    \caption{So sánh năng lực hạ tầng giữa các nhà cung cấp Cloud}
    \label{tab:cloud_comparison}
    \small
\end{table}
\\ 
\noindent Nhóm chọn kết hợp Bizfly Cloud để vận hành các tính năng hằng ngày nhằm giảm chi phí và tăng tốc độ, đồng thời dùng Amazon (AWS) để lưu trữ dữ liệu quan trọng nhằm đảm bảo tính an toàn tuyệt đối.         
\\ \\ \\
\noindent Dựa trên sự kết hợp giữa các nhà cung cấp, chi phí phần cứng để vận hành hệ thống được cụ thể hóa như sau
\begin{table}[htbp]
    \centering
    \begin{tabularx}{\textwidth}{@{}l l r r@{}}
        \toprule
        \textbf{Hạng mục} & \textbf{Nhà cung cấp} & \textbf{Đơn giá/Tháng} & \textbf{Tổng 12 tháng (VNĐ)} \\
        \midrule
        Tính toán & Bizfly Cloud (Basic) & 300.000 & 3.600.000 \\
        Quản lý dữ liệu & AWS & 450.000 & 5.400.000 \\
        Lưu trữ \& CDN & AWS S3 \& Cloudflare & 150.000 & 1.800.000 \\
        \midrule
        \multicolumn{3}{@{}l}{\textbf{Tổng chi phí dự kiến}} & \textbf{10.800.000} \\
        \bottomrule
    \end{tabularx}
    \caption{Dự toán chi tiết hạ tầng Cloud năm thứ nhất}
    \label{tab:cloud_budget_year1}
    \small
\end{table}
\\ \\ \\
Để duy trì dịch vụ cho 80 cụm sân, dự toán chi phí hàng tháng được phân bổ như sau:
\begin{table}[H]
    \centering
    \begin{tabularx}{\textwidth}{@{}l X r@{}}
        \toprule
        \textbf{Hạng mục} & \textbf{Mô tả chi tiết} & \textbf{Dự toán (VNĐ)} \\
        \midrule
        Hệ thống lưu trữ \& Kết nối & Duy trì máy chủ, cơ sở dữ liệu khách hàng, bản đồ số và hệ thống gửi thông báo/SMS cho chủ sân & 2.500.000 \\
        \addlinespace
        Tiếp thị \& Phát triển thị trường & Chi phí tìm kiếm khách hàng (chủ sân), quảng cáo trên mạng xã hội và các chương trình khuyến khích giới thiệu người dùng mới & 10.000.000 \\
        \addlinespace
        Quản lý \& Hỗ trợ kỹ thuật & Chi phí nhân sự trực tuyến để giải quyết các vấn đề kỹ thuật phát sinh và hỗ trợ chủ sân vận hành phần mềm & 15.000.000 \\
        \midrule
        \textbf{Tổng chi phí} & & \textbf{27.500.000} \\
        \bottomrule
    \end{tabularx}
    \caption{Dự toán chi phí vận hành hàng tháng (OPEX)}
    \label{tab:opex_estimate}
    \small
\end{table}

\subsubsection{Mô hình doanh thu và Điểm hòa vốn}
Nền tảng vận hành theo mô hình cung cấp phần mềm dưới dạng dịch vụ thuê bao định kỳ (SaaS). Chủ sân không cần trả hoa hồng trên mỗi lượt khách đặt, thay vào đó sẽ trả phí duy trì cố định để dễ dàng quản lý lợi nhuận. \\ 

\noindent Các gói dịch vụ dành cho chủ sân: được thiết kế phù hợp với quy mô kinh doanh của từng cụm sân
\begin{itemize}
    \item \textbf{Gói Tiêu chuẩn (Standard)}: dành cho các cụm dưới 5 sân, mức phí 299.000 VNĐ/tháng.
    \item \textbf{Gói Doanh nghiệp (Enterprise)}: dành cho các cụm trên 5 sân, mức phí 499.000 VNĐ/tháng.
\end{itemize} 
\vspace{0.5cm}
\noindent Ước tính doanh thu hàng tháng khi hệ thống đạt quy mô 80 cụm sân (giả định 60\% gói Enterprise và 40\% gói Standard)
\begin{itemize}
    \item \textbf{Doanh thu từ phí dịch vụ (SaaS)}: uớc tính đạt 33.520.000 VNĐ/tháng (48 $\times$ 499.000 + 32 $\times$ 299.000)
    \item \textbf{Doanh thu từ quảng cáo}: uớc tính 5.000.000 VNĐ/tháng từ việc hiển thị thương hiệu của các nhãn hàng thể thao tại các vị trí ưu tiên trên hệ thống.
    \item \textbf{Lợi nhuận ròng}: sau khi trừ mọi chi phí, hệ thống mang về lợi nhuận khoảng 11.020.000 VNĐ/tháng.
    \item \textbf{Thời gian hoàn vốn}: với tổng vốn đầu tư ban đầu ước tính 150 - 200 triệu VNĐ, dự án dự kiến đạt điểm hòa vốn sau 14 - 18 tháng hoạt động ổn định.
\end{itemize}


\subsection{Cơ cấu tổ chức và Mô tả công việc}
\subsubsection{Sơ đồ tổ chức theo giai đoạn phát triển}
Cơ cấu nhân lực của CourtMate được chia làm 2 giai đoạn để tối ưu hóa chi phí nhân sự:
\begin{enumerate}
    \item \textbf{Giai đoạn Xây dựng (Startup/MVP)}: tập trung 100\% vào đội ngũ kỹ sư để hoàn thiện sản phẩm.
    \item \textbf{Giai đoạn Vận hành (Scale-up)}: giữ nguyên đội ngũ kỹ thuật và bổ sung bộ phận kinh doanh (Sales) để trực tiếp tiếp cận các chủ sân.
\end{enumerate}

\subsubsection{Mô tả công việc các vị trí cốt lõi}
Để đảm bảo tính chuyên môn hóa, các thành viên được phân bổ nhiệm vụ theo hướng hỗ trợ mô hình kinh doanh:
\begin{itemize}
    \item \textbf{Hệ thống lõi (Backend Lead - Trần Thanh Phúc)}: thiết kế và vận hành logic của hệ thống, đảm bảo quy trình đặt sân, hủy sân và báo cáo tài chính diễn ra chính xác.
    \item \textbf{Trải nghiệm người dùng (Frontend Lead - Trần Minh Quân)}: phụ trách phát triển về giao diện quản lý lịch sân (kéo-thả) giúp chủ sân thao tác dễ dàng và theo dõi các biểu đồ kinh doanh trực quan.
    \item \textbf{Giao diện khách hàng (Frontend Developer - Nguyễn Đăng Khoa)}: xây dựng các tính năng tìm kiếm sân trên bản đồ, quy trình đăng ký thông tin sân và hệ thống quét mã QR để khách hàng check-in tại sân
    \item \textbf{Bảo mật và Tích hợp (Security \& Integration Engineer - Phạm Văn Quốc Việt)}: đảm bảo an toàn cho tài khoản người dùng, xử lý các giao dịch thanh toán trực tuyến và tích hợp hệ thống hóa đơn điện tử.
    \item \textbf{Kiểm soát Chất lượng \& Vận hành máy chủ (QA/QC \& DevOps - Lê Võ Đăng Khoa)}: kiểm tra lỗi hệ thống trước khi ra mắt, đảm bảo máy chủ luôn hoạt động 24/7 và có khả năng chịu tải cao.
    \item \textbf{Phát triển tính năng bổ trợ (Backend Developer - Nguyễn Trần Bảo Khoa)}: xử lý các nghiệp vụ liên quan đến việc giữ chỗ (lock slot) để tránh tình trạng hai khách hàng đặt trùng một giờ sân, và tối ưu hóa công cụ tìm kiếm sân gần nhất.
\end{itemize}






